\noindent326\hfill\textbf{BRIEFWECHSEL.\quad ZUR NETZAUSGLEICHUNG.}
\begin{center}\textsc{Gauss} an \textsc{Gerling}.\quad Göttingen, 14. November 1838.\end{center}

.\ .\ .\ .\ .\ .\ Rücksichtlich der Formeln für die Anzahl der Bedingungsgleichungen
habe ich gar nichts gegen die Bekanntmachung. Allein es ist mir
\textls[120]{jetzt} unmöglich, die Ausnahmsfälle auf eine präcise und erschöpfende Art
zu bezeichnen, und ich weiss auch nicht, ob dies überhaupt \textls[120]{möglich} sein
wird, ohne tiefer eingreifende Entwickelungen damit zu verbinden. S[ine] m[editatione]
dürfte es zureichen zu sagen, dass wenn in einem System von $p$
Punkten, zwischen denen es $l$ Verbindungslinien gibt, die Richtung jeder Verbindungslinie
an \textls[120]{beiden} Endpunkten gemessen, d.\ i.\ an eine andere Richtungslinie
des Systems geknüpft ist, es

$l-p+1$ von einander unabhängige Bedingungsgleichungen aus blossen
Winkelsummen und ausserdem noch

$l-2p+3$ \textls[120]{andere} Bedingungsgleichungen geben muss, die in den
\textls[120]{gewöhnlichsten} Fällen sämmtlich unter die Kategorie der in dem Supplem.
Theor.\ p.\ 31 unter III[\textsuperscript{*})] angeführten fallen.
\par\medskip
\noindent\rule{.24\textwidth}{.4pt}\par\smallskip
{\footnotesize\noindent\hspace*{1.5em}[\textsuperscript{*}) Band IV, S.\ 81.]\par}
\vspace{2em}
\begin{center}\rule{.18\textwidth}{.4pt}\end{center}
